\section{Finite-angle distinguishability on unitary orbits}\label{sec:orbit84}
Fix an ordered measurement $E$ and a Hermitian generator $G$, and write
\begin{equation}\label{eq:unitorbit84}
 E_j(t)=e^{itG}E_je^{-itG},\qquad H_j=i[G,E_j].
\end{equation}
Both $E$ and $G$ are known in this pair-comparison problem.  A tester
must use the same controls under the two hypotheses $E$ and $E(t)$;
its design may depend on this known pair.  It has access to the
classical device labels and its own reference and memory, never to
an unobserved measurement environment.

The Hamiltonian-not-in-Kraus-span criterion and its quantum-error-
correction interpretation are established in quantum channel metrology
\cite[Theorems~1--3]{ZJ84}.  The support criterion below is its
measurement-channel specialization.  We prove a finite-angle,
finite-use trace-distance statement directly, including an explicit
accumulated error estimate.  Thus no lower bound is inferred solely
from an infinitesimal Fisher-information scaling.

\begin{theorem}[Finite-use trichotomy on every fixed unitary orbit]\label{thm:orbittrichotomy84}
Let $d,N\ge1$, $k\ge2$, and let $E$ and $G$ be as in
\eqref{eq:unitorbit84}, with no rank or strict-positivity assumption.
The following alternatives are exhaustive.
\begin{enumerate}
\item If $[G,E_j]=0$ for every $j$, then $D_N(E,E(t))=0$ for all
$N$ and all $t$.
\item If $H\ne0$ and $G\in\mathcal W_E$, there are constants
$c,C,t_0>0$, depending on $E,G$, such that
\begin{equation}\label{eq:orbitsql84}
 c\min\{1,\sqrt N|t|\}\le D_N(E,E(t))
 \le C\min\{1,\sqrt N|t|\},\qquad |t|\le t_0.
\end{equation}
The lower bound is attained by a product-input, nonadaptive experiment.
\item If $G\notin\mathcal W_E$, there are constants
$c,C,t_0>0$, depending on $E,G$, such that
\begin{equation}\label{eq:orbithl84}
 c\min\{1,N|t|\}\le D_N(E,E(t))
 \le C\min\{1,N|t|\},\qquad |t|\le t_0.
\end{equation}
The lower bound admits a common adaptive experiment carrying a logical
qubit between calls, with a reference of dimension at most $2d$ at
an individual call and ideal trusted controls.
\end{enumerate}
The constants are not uniform as the measurement or generator changes.
All upper bounds allow arbitrary retained reference and adaptive memory,
and bounded public stopping.  This theorem classifies a fixed unitary
orbit; it does not assert two-sided equivalence of the midpoint
covariance modulus on arbitrary pairs in the complete measurement body.
\end{theorem}

We first record the error-correction statement in the form needed for
finite angles.  Its exact-correction premise is the usual
Knill--Laflamme condition \cite{KL84}; a proof is included to specify
which systems are available to the recovery.

\begin{lemma}[A corrected logical channel and its quadratic error]\label{lem:correctedchannel84}
Let $\mathcal M$ be a channel on an input space $\mathsf A$, and
let $J:\C^2\longrightarrow\mathsf A\otimes\mathsf R$ be an isometry.
Suppose a channel $\mathcal R$ from the actual output of
$\mathcal M\otimes\operatorname{Id}_{\mathsf R}$ to $\C^2$ satisfies
\[
 \mathcal R\circ(\mathcal M\otimes\operatorname{Id}_{\mathsf R})
          \circ\mathcal J=\operatorname{Id}_2,
 \qquad \mathcal J(X)=JXJ^*.
\]
For a Hermitian $G$ on $\mathsf A$, put
$G_L=J^*(G\otimes I)J$, $\mathcal U_t(X)=e^{-itG_L}Xe^{itG_L}$,
and
\[
 \Phi_t=\mathcal R\circ(\mathcal M\otimes\operatorname{Id}_{\mathsf R})
       \circ\operatorname{Ad}(e^{-it(G\otimes I)})\circ\mathcal J.
\]
In unhalved diamond norm,
\begin{equation}\label{eq:logicalerror84}
 \|\Phi_t-\mathcal U_t\|_\diamond
 \le4t^2\|G\|_{\mathrm{op}}^2,
 \qquad
 \|\Phi_t^{\circ m}-\mathcal U_t^{\circ m}\|_\diamond
 \le4mt^2\|G\|_{\mathrm{op}}^2.
\end{equation}
\end{lemma}
\begin{proof}
Let $\mathcal T=\mathcal R\circ(\mathcal M\otimes\operatorname{Id})$
and choose Kraus operators $V_\ell$ for $\mathcal T$.  Since
$\mathcal T\circ\mathcal J$ is the identity channel, whose Choi
operator has rank one, $V_\ell J=c_\ell I_2$ with
$\sum_\ell|c_\ell|^2=1$.  Trace preservation on the entire input
space gives
\[
 \sum_\ell c_\ell^*V_\ell
 =\sum_\ell(V_\ell J)^*V_\ell
 =J^*\sum_\ell V_\ell^*V_\ell=J^*.
\]
It follows by differentiating that $\Phi_0=\operatorname{Id}_2$
and $\Phi'_0=-i[G_L,\,\cdot\,]$.
For any Hermitian $B$, the generator of $\operatorname{Ad}(e^{-itB})$
has diamond norm at most $2\|B\|_{\mathrm{op}}$.  Its second
 derivative has diamond norm at most $4\|B\|_{\mathrm{op}}^2$,
since unitary channels have diamond norm one.  Taylor's integral
remainder and contractivity under the surrounding channels bound
the remainder for $\Phi_t$ by $2t^2\|G\|_{\mathrm{op}}^2$.
The remainder for $\mathcal U_t$ has the same bound because
$\|G_L\|_{\mathrm{op}}\le\|G\|_{\mathrm{op}}$.  Their zeroth
and first derivatives agree, proving the first inequality.
Telescoping the two $m$-fold compositions of channels proves the second.
\end{proof}

\begin{lemma}[A reference-preserving code from the support complement]\label{lem:supportcode84}
Suppose $G\notin\mathcal W_E$ and put
\begin{equation}\label{eq:signstates84}
 A=G-\Pi_{\mathcal W_E}G,\qquad
 s=\tr A_+=\tr A_->0,\qquad \rho_\pm=A_\pm/s.
\end{equation}
There is a code $J:\C^2\to\C^d\otimes\mathsf R$ with
$\dim\mathsf R\le2d$, corrected exactly by a channel acting only
on the classical label and $\mathsf R$ after a call to $E$.
The logical generator $G_L=J^*(G\otimes I)J$ is diagonal with gap
\begin{equation}\label{eq:logicalgap84}
 \Delta=\tr\bigl(G(\rho_+-\rho_-)\bigr)
        =\|A\|_{\mathrm{HS}}^2/s>0.
\end{equation}
\end{lemma}
\begin{proof}
The identity belongs to $\mathcal W_E$ because $\sum_jE_j=I$;
thus $A$ is traceless.  Its nonzero positive and negative parts have
equal trace.  Also
\[
 P_j\rho_+P_j=P_j\rho_-P_j\quad\text{for every }j.
\]
Purify $\rho_+$ and $\rho_-$ in $\C^d\otimes\C^d$ and attach
orthogonal flags to the second factor.  Map the two logical basis
vectors to these flagged purifications.  They are orthonormal,
and the flags make all off-diagonal code matrix elements of
$B\otimes I$ vanish.  The two diagonal elements are equal for
every $B\in P_jM_d(\C)P_j$ by the compression identity above.

Choose measurement Kraus operators
$L_{j,a}=|j\rangle\langle a|\sqrt{E_j}$, $1\le a\le d$.
Products with different labels are zero; for one label their span
is $P_jM_d(\C)P_j$.  Hence, with $\mu=(j,a)$,
\begin{equation}\label{eq:kl84}
 J^*(L_\mu^*L_\nu\otimes I)J=\lambda_{\mu\nu}I_2
\end{equation}
for a positive matrix $\lambda$ of trace one.  Diagonalize
$\lambda$ by a unitary change of Kraus representation.  If its
positive eigenvalues are $\lambda_b$, the corresponding operators
$(\widetilde L_b\otimes I)J/\sqrt{\lambda_b}$ are isometries
with mutually orthogonal ranges.  A channel projects onto those
ranges and applies their inverse isometries; on the orthogonal
complement it prepares a fixed logical state.  This is completely
positive and trace preserving on the entire output space, and it
corrects the code exactly.  Its input space is just the actual
classical-label register tensor $\mathsf R$.  It does not contain
the lost input or a dilation environment.  Equivalently, since the
actual state is label diagonal, the same channel may first read the
label and then act on the retained reference.

The logical generator is diagonal by the flags, with entries
$\tr(G\rho_\pm)$.  Orthogonality of $A$ to
$\Pi_{\mathcal W_E}G$ gives
$\tr(GA)=\tr(A^2)$, proving \eqref{eq:logicalgap84}.
\end{proof}

\begin{proof}[Proof of Theorem~\ref{thm:orbittrichotomy84}]
The first assertion follows directly from \eqref{eq:unitorbit84}.
For the second, Proposition~\ref{prop:unitaryrange84} gives
\[
 h^2=\langle H,\mathcal C_E^\dagger H\rangle<\infty.
\]
Conjugating all tuple components by $e^{itG}$ intertwines the
covariances and their tangent directions, so this value is constant
along the orbit.  Inverse spectral comparison gives
$g_{N,E(t)}(E'(t))\le\sqrt N h$.  The complete-body path theorem
therefore yields
\begin{equation}\label{eq:orbitupper84}
 D_N(E,E(t))\le\min\{2,\sqrt{k+1}\,h\sqrt N|t|\}.
\end{equation}
There is an index $j$ with $H_j\ne0$.  Choose a unit vector $v$
with $\langle v,H_jv\rangle\ne0$.  The two Bernoulli parameters
obtained by input $v$ and the event ``label $j$'' have gap at least
$c_0|t|$ on a sufficiently small two-sided interval.  The finite
Bernoulli-product bound used in
Theorem~\ref{thm:multioutcomemetric82} gives
$(1/128)\min\{1,c_0\sqrt N|t|\}$.  This proves
\eqref{eq:orbitsql84}, including all finite $N$.

For the third alternative, use Lemma~\ref{lem:supportcode84}.
The equality
$\mathcal M_{E(t)}=\mathcal M_E\circ\operatorname{Ad}(e^{-itG})$
holds also with a retained reference.  Encode, make one device call,
recover, and repeat $m\le N$ times.  Under $E$ the logical channel
is the identity.  Under $E(t)$ it is $\Phi_t^{\circ m}$ from
Lemma~\ref{lem:correctedchannel84}.  Start in the equal superposition
of the two logical basis states.  The ideal logical unitary produces
unhalved trace distance $2|\sin(mt\Delta/2)|$ from that state.
Consequently, for every real $t$,
\begin{equation}\label{eq:finiteanglelower84}
 D_N(E,E(t))\ge
 \max_{1\le m\le N}
 \left(2|\sin(mt\Delta/2)|-4mt^2\|G\|_{\mathrm{op}}^2\right)_+.
\end{equation}
This bound concerns one specified common tester for each $m$,
not an inaccessible choice of a different recovery under the two hypotheses.

For $t\ne0$ choose
\[
 |t|\le t_0:=\min\left\{\frac1{2\Delta},
                       \frac{\Delta}{4\pi\|G\|_{\mathrm{op}}^2}\right\},
 \qquad
 m=\min\left\{N,\left\lfloor\frac1{|t|\Delta}\right\rfloor\right\}.
\]
Then $x=m|t|\Delta\le1$ and
$x\ge\tfrac12\min\{1,N|t|\Delta\}$.
Since $2\sin(x/2)\ge2x/\pi$ and the error in
\eqref{eq:finiteanglelower84} is at most $x/\pi$, we obtain
\begin{equation}\label{eq:explicitlocalhl84}
 D_N(E,E(t))\ge\frac1{2\pi}\min\{1,N|t|\Delta\}.
\end{equation}
The elementary hybrid estimate for the input unitaries gives
$D_N(E,E(t))\le\min\{2,2N|t|\|G\|_{\mathrm{op}}\}$,
which proves \eqref{eq:orbithl84}.  The zero-angle case is immediate.
The protocol uses deterministic $m\le N$ calls, hence belongs to
the class allowing bounded public stopping.  The upper estimates
hold for that entire class.  Finally if the first alternative held
with $G\notin\mathcal W_E$, then $H=0$ would contradict
Proposition~\ref{prop:unitaryrange84}; the alternatives are disjoint.
\end{proof}

\begin{corollary}[Full-rank, rank-one and projective measurements]\label{cor:orbitexamples84}
If one effect is positive definite, every nonconstant unitary orbit
has the square-root law \eqref{eq:orbitsql84}.  If all nonzero
effects have rank one, every nonconstant orbit has that law precisely
when the measurement is informationally complete, meaning that its
effects span $\mathcal H_d$.  If such a rank-one measurement is not
informationally complete, some generator has the linear law
\eqref{eq:orbithl84}.  For a projection-valued measurement every
orbit is either constant or has the linear law.
\end{corollary}
\begin{proof}
A full-rank support contributes $\mathcal H_d$ to $\mathcal W_E$.
For rank-one supports, $P_j\mathcal H_dP_j=\R E_j$, so that
$\mathcal W_E$ is exactly the effect span.  If this span is proper,
a nonzero element of its orthogonal complement is a generator
outside $\mathcal W_E$ and cannot commute with all effects by
Proposition~\ref{prop:unitaryrange84}.  For a projection-valued
measurement $\mathcal W_E$ is the block-diagonal Hermitian space
of the orthogonal supports, which is exactly the space of generators
commuting with every effect.  Apply Theorem~\ref{thm:orbittrichotomy84}.
\end{proof}

\begin{proposition}[Support spans under output processing]\label{prop:supportprocessing84}
Let $T_{aj}\ge0$ satisfy $\sum_aT_{aj}=1$, and put
$F_a=\sum_jT_{aj}E_j$.  Then $\mathcal W_E\subseteq\mathcal W_F$.
Thus classical processing of the ordered output cannot turn a
nonconstant square-root orbit into a linear orbit.  A reversible
splitting of labels preserves the support span and the classification.
\end{proposition}
\begin{proof}
For positive matrices, the kernel of a positive weighted sum is the
intersection of the kernels of its positively weighted summands.
Thus the support of $F_a$ contains $P_j$ whenever $T_{aj}>0$.
Every column has such a row, giving
$P_j\mathcal H_dP_j\subseteq\mathcal W_F$.
This proves the inclusion.  The orbit commutes with classical
processing, so Theorem~\ref{thm:orbittrichotomy84} applies with the
same generator.  A reverse stochastic map gives the opposite
inclusion for a reversible splitting.
\end{proof}

\begin{remark}[What the lower construction does not assert]\label{rem:orbitresources84}
The code and recovery depend on the known pair through $E,G$.
They do not constitute a common estimator of an unknown measurement.
Their existence is proved in finite-dimensional quantum mechanics;
polynomial gate synthesis, physical execution, and parameter-uniform
bounds for approximate controls are not inferred.  Formula
\eqref{eq:finiteanglelower84} quantifies the finite-angle channel
remainder under ideal controls.  It is not a claim of perfect
finite-use discrimination.  The full-body entropy and growing-$k$
learning questions are separate from this orbit classification.
\end{remark}
